\documentclass[12pt,a4paper]{report}
\usepackage[pdftex]{color,graphicx}
\usepackage[brazil]{babel}
\usepackage[utf8]{inputenc}
\usepackage[T1]{fontenc}
\usepackage[a4paper,left=1.18in,right=1.18in,top=0.98in,bottom=0.98in]{geometry}
\usepackage{fancyhdr}

\hyphenation{re-a-li-za-dos pre-di-to mo-de-la-gem mo-di-fi-ca-ções re-fe-rên-cia res-pec-ti-va-men-te pa-ra-mag-né-ti-co di-fe-ren-te li-mi-te}


\begin{document}
\fontsize{15pt}{17pt}\selectfont

\begin{titlepage}
\begin{center}
\includegraphics[width=0.2\textwidth]{./logo1.pdf}\\[0.2cm]
\textsc{UNIVERSIDADE FEDERAL DO PIAUÍ}\\[0.2cm]
\textsc{CENTRO DE CIÊNCIAS DA NATUREZA}\\[0.2cm]
\textsc{DEPARTAMENTO DE FÍSICA}\\[2.0cm]
{\textbf{Programa PIBIC-UFPI (2014/2015) – Relatório Parcial}}\\[1.0cm]
{\textbf{PLANO DE TRABALHO}}\\[0.4cm]
{\textbf{``Estudo teórico do efeito de pressões extremas em nanotubos de
carbono''}}\\[1.5cm]
{\textbf{PROJETO DE PESQUISA}}\\[0.4cm]
{\textbf{``ESTUDO COMPUTACIONAL DAS PROPRIEDADES VIBRACIONAIS E ELETRÔNICAS 
SISTEMAS HÍBRIDOS BASEADOS EM NANOTUBOS DE CARBONO SUJEITOS A DEFORMAÇÕES ESTRUTURAIS''}}\\[2.0cm]

{Orientando: Seu nome aqui}\\[1.0cm]
{Orientador: Acrisio Lins de Aguiar}\\[1.0cm]
{Teresina--PI, Fevereiro de 2015}
\end{center}

\end{titlepage}
\setcounter{page}{2}

\chapter*{INTRODUÇÃO}

Este trabalho teve como objetivo simular diferentes tipos de Nanotubo de Carbono computacionalmente. 
Foi utilizada a linguagem de programação C++  e o compilador dev-C++. Quanto a visualização do resultado final foi usado o software de visualização
Chemcraf. Muito antes do termo nano ser adotado na comunidade cientifica, o fisico  norte-americano Richard Phillips Feynman (1918-1988) 
em sua palestra “There’s Plenty of Room at the Bottom” (Há muito espaço lá em baixo) ocorrida em 1959 num encontro da 
Sociedade Americana de Fisica intrigou varios ouvintes ao falar em manipulação de objetos em escala atômica. 
A partir de 1987, foram surgindo os primeiros trabalhos relacionados a nanotecnologia. Foi com o invento do 
microscopio de tunelamento (STM- Scanning Tunneling Microscope) que se tornou possivel a manipulacão de materiais nanométricos.

A Nanociência trata do estudo da manipulação  de materiais na ordem do namometro, na ordem de  \(10^{-9}\) metros, 
afim de se criar e modificar estruturas maiores.Conhecimentos dela são aplicadas em diversas áreas como a Medicina, 
Engenharias, dentre  outras ciencias.
O carbono,unico elemento da tabela periodica capaz de formar alotropos de diferentes dimensionalidades,  
é um elemento quimico de grande utildade. Elemento de maior abundancia no universo, ele é dito como o elemento 
base pra vida terrestre. Também importante  no meio cientifico, o carbono é comumente encontrado na natureza na 
forma de seus alotropos, dentre eles o grafite e o diamante. O diamante é umas das mais fortes estruturas já 
encontradas mas é do grafite que a folha de grafeno é obtida. E é esse o material feito de atomos de carbono que tem sido alvo de grandes projetos.

O grafeno nada mais é que uma estrutura bidimensional feita de atomos de carbonos com a espessura do mesmo.  
A partir do grafeno obtemos o Nanotubo de Carbono.
Nanotubos de carbono (CNTs), também alotropo do carbono,  são estruturas cilíndricas  e ocas geradas pelo 
enrolar de uma folha de grafeno, tendo assim a superfície muito fina,  espessura de um átomo de carbono. 
Podem ser separados em dois conjuntos:  de uma única só superfície de grafeno (SWNTs) ou mais folhas enroladas 
de forma concêntricas (MWNTs).  CNTs são altamente resistentes à ruptura e possuem um grande poder relacionado 
a condutibilidade térmica. As propriedades estruturais do Nanotubo de Carbono definem suas características 
físicas e suas aplicações especificas, podendo ser ótimos condutores elétricos.
Há controversias sobre a descoberta do Nanotubo de Carbono, entretanto foi a parti do trobalho do Sumio Iijima, 
um físico japonês, feito em 1991 que os Nanotubos de Carbono passaram a ter destaque no meio cientifico.

Em 1985 houve a descoberta de outro alotropo do carbono, o fulereno. Sendo uma  nanomolecula esferoidal 
e ao contrario dos outros alotropos esse possui uma forma molecular.
Com tantas descobertas  envolvendo os alotropos do carbono parte consideravel do meio cientifico se voltou 
para o estudo de suas propriedades e utilidades.



\chapter*{REVISÃO DA LITERATURA}

\chapter*{METODOLOGIA}

\chapter*{RESULTADOS E DISCUSSÃO}

\chapter*{CONCLUSÃO}

\chapter*{REFERÊNCIAS}
\bibliographystyle{edu_style2}
\bibliography{bibs}


\end{document}
